% Document class and font size
\documentclass[a4paper,9pt]{extarticle}

% Packages
\usepackage[utf8]{inputenc} % For input encoding
\usepackage{geometry} % For page margins
\geometry{letterpaper, margin=0.75in} % Set paper size and margins
\usepackage{titlesec} % For section title formatting
\usepackage{enumitem} % For itemized list formatting
\usepackage{hyperref} % For hyperlinks
\usepackage{tabularx}
\usepackage{fancyhdr}

% Formatting
\setlist{noitemsep} % Removes item separation
\titleformat{\section}{\large\bfseries}{\thesection}{1em}{}[\titlerule] % Section title format
\titlespacing*{\section}{0pt}{\baselineskip}{\baselineskip} % Section title spacing
%%%%%%%%%%

% Begin document
\begin{document}

% Disable page numbers
\pagestyle{fancy}
\renewcommand{\headrulewidth}{0pt}
\fancyhead{}
\fancyhead[L]{\textit{Baibhab Karmakar}}
\fancyhead[R]{\textit{August 2026}}
\thispagestyle{empty} % Remove header from the first page

% Header
\begin{flushleft}
\textbf{\LARGE Baibhab Karmakar}\\[2pt] % Name
Master's Student in Earth Sciences, IISER Kolkata, Kolkata, India
\\ \href{mailto:bk22ms064@iiserkol.ac.in}{bk22ms064@iiserkol.ac.in} | \href{tel:+918250734897}{+91 8250734897} | \href{https://www.linkedin.com/in/baibhab-karmakar-237249195/}{LinkedIn} % Contact info
\end{flushleft}

% Research Interests
\section*{RESEARCH INTERESTS}
\noindent
Over the past two years, I have worked on earthquake source characterization, slip inversion, surface-wave amplification, and machine learning-based earthquake relocation, using observational and computational approaches to analyse seismic data.  I aim to deepen my expertise through PhD research in understanding earthquake processes, seismic wave propagation, and Earth structure, with a particular interest in computational geophysics and data-driven methods in seismology.
% Education Section
\section*{EDUCATION}
\noindent
\textbf{Indian Institute of Science Education and Research (IISER) Kolkata}, Kolkata, India \hfill 2026 -- Present\\
Master's Student in Earth Sciences \\
Thesis Title: Earthquake Relocation Using Machine Learning % Additional info

\noindent\\
\textbf{Indian Institute of Science Education and Research (IISER) Kolkata}, Kolkata, India \hfill 2022 -- 2026\\
BS-MS, Earth Sciences Major and Computer Science Minor \hfill CGPA: 8.09/10.00

\noindent\\
\textbf{Higher Secondary Examination} \hfill 2022\\
Score: 95.8\%

\noindent\\
\textbf{Madhyamik Examination} \hfill 2020\\
Score: 96.29\%

% Research Experience Section
\section*{RESEARCH EXPERIENCE}
\noindent
\textbf{Computational Seismology Laboratory, IISER Kolkata} \hfill Kolkata, India\\
\textit{Master's Researcher, under Prof. Supriyo Mitra} \hfill 2026 | Present
\begin{itemize}
    \item Applying QSeek, a machine-learning-based earthquake detection and localization framework, to continuous seismic data from the Kashmir Himalaya.
    \item Combining neural-network phase picking (via SeisBench pickers such as PhaseNet/EQTransformer) with stacking-migration and adaptive octree localization to build a data-driven earthquake catalog for the region.
    \item Benchmarking detection completeness and location accuracy of the QSeek-derived catalog against the existing regional catalog to characterize seismicity patterns in the Kashmir Himalaya. 
\end{itemize}
\noindent
\textbf{Institut de Physique du Globe de Paris (IPGP)} \hfill Paris, France\\
\textit{Summer Research Intern, under Dr. Martin Vall\'ee} \hfill June 2026 | July 2026
\begin{itemize}
    \item Performed joint kinematic slip inversions for the 2015 and 2025 Charlie-Gibbs Transform Fault earthquakes using teleseismic $P$/$SH$ waveforms and moment-constrained Empirical Green's Function relative source time functions (RSTFs), resolving unilateral westward rupture propagation in both events.
    \item Used RBF interpolation and polygon overlap analysis to quantify asperity overlap and disjointness  between the two ruptures at different slip thresholds.
    \item Revealed contrasting rupture patterns on the Charlie-Gibbs fault, with the 2015 earthquake dominated by a localized high-slip asperity, while the 2025 event exhibited more distributed slip and limited re-slip at the 2015 peak-slip region.
    \item Improved the seismic-data fit by incorporating spatially variable fault dip as an inversion parameter, while noting that the resulting dip distribution lacks independent observational constraints.
\end{itemize}

\noindent
\textbf{IISER Pune} \hfill Pune, India\\
\textit{Research Intern} \hfill 2025
\begin{itemize}
    \item Investigated surface-wave amplification across the Shillong Plateau–Bengal Basin boundary using 1D shear-wave velocity profiles, focusing on the effects of low shear velocities and impedance contrasts at the basin edge.
    \item Modeled surface-wave propagation using SWRT (semi-analytical solutions across transversely aligned vertical discontinuities, incorporating higher-mode conversion) alongside SPECFEM2D full-waveform simulations; both methods predicted 2--5$\times$ higher amplitudes in the basin than in crystalline terrain over 0.01--0.15 Hz.
    \item Validated model predictions against a regional earthquake recorded at both sites, finding good agreement on the transverse component; presented as a poster at AGU Annual Meeting 2025.
\end{itemize}
\noindent
\textbf{Environmental Science Laboratory, IISER Kolkata} \hfill Kolkata, India\\
\textit{Research Assistant} \hfill 2024
\begin{itemize}
    \item Assisted in environmental data collection and analysis on pesticide contamination in milk samples.
\end{itemize}

% Teaching Experience Section
\section*{TEACHING EXPERIENCE}
\noindent
\textbf{IISER Kolkata} \hfill Kolkata, India\\
\textit{Teaching Assistant - Geophysics and Hydrology} \hfill Autumn 2026 - Present
\begin{itemize}
    \item Assisting in tutorials and practical sessions for undergraduate students.
    \item Supporting students in conceptual understanding and problem-solving in geophysics and hydrology.
\end{itemize}

\noindent
\textbf{IISER Kolkata} \hfill Kolkata, India\\
\textit{Teaching Assistant - Introduction to Earth Science} \hfill Autumn 2025
\begin{itemize}
    \item Assisted in conducting tutorials and practical sessions for undergraduate students.
    \item Supported grading, doubt-clearing sessions, and conceptual understanding of core Earth Science topics.
\end{itemize}

% Projects Section
\section*{PROJECTS}
\noindent
\textbf{Receiver Function Tool} \hfill IISER Kolkata, India\\
\textit{Self-directed Computational Project}
   \hfill 2025

\begin{itemize}
     \item Built an automated and interactive tool for crustal thickness estimation using receiver-function analysis.
     \item Developed workflows for seismic interpretation and theoretical travel-time modelling.
\end{itemize}

\section*{PUBLICATIONS}
\noindent
\textbf{Journal paper}\\
\textup{In Preparation}
\begin{itemize}
    \item Karmakar, B., Vall\'ee, M.. Ten-year fault reactivation on the Charlie-Gibbs Transform Fault: joint kinematic inversion and slip asperity distribution analysis of the 2015 and 2025 earthquakes (working title, in preparation).
\end{itemize}

\section*{PRESENTATIONS}
\noindent
\textbf{Poster Presentation - AGU Annual Meeting 2025}\\
\textit{American Geophysical Union (AGU)}
\begin{itemize}
    \item Presented my work on seismic surface wave amplification in the Bengal Basin.
    \item Discussed observational and modelling results with an international research audience.
\end{itemize}
\section*{SELECTED COURSES}
\begin{tabularx}{1\textwidth}{
    >{\raggedright\arraybackslash}X
    >{\raggedright\arraybackslash}X }
\textbf{Master's Courses} & \textbf{Bachelor's Courses} \\[3pt]
\begin{minipage}[t]{\linewidth}
\begin{itemize}
    \item Master's Thesis Research in Computational Seismology
\end{itemize}
\vspace{6pt}
\textbf{Additional Courses}
\begin{itemize}
    \item Data Science and Machine Learning
    \item Artificial Intelligence for Data Science
\end{itemize}
\end{minipage}
&
\begin{minipage}[t]{\linewidth}
\begin{itemize}
    \item Seismology
    \item Seismology Lab
    \item Geophysics
    \item Waves and Optics
    \item Structural Geology
    \item Geodynamics
\end{itemize}
\end{minipage}
\end{tabularx}
\section*{AWARDS \& FELLOWSHIPS}
\textbf{Innovation in Science Pursuit for Inspired Research (INSPIRE)}  \hfill India\\ 
Department of Science and Technology, Government of India \hfill 2022 -- Present\\ \\
\textbf{JBNSTS Junior Fellow}  \hfill India\\ 
Jagadis Bose National Science Talent Search \hfill 2020 -- 2022
\section*{SKILLS}
\begin{itemize}
    \item \textbf{Programming:} Python, C, MATLAB, Shell scripting
    \item \textbf{Python Libraries:} NumPy, Pandas, SciPy, Matplotlib, ObsPy, SeisBench, Sklearn, TensorFlow and etc.
    \item \textbf{Geophysical Software:} SAC, SPECFEM2D, SWRT, ASPECT(basic), Pyrocko
    \item \textbf{Machine Learning:} Supervised and unsupervised learning, data-driven modelling
    \item \textbf{Other:} Data Structures and Algorithms, Git, LaTeX, Microsoft Office Suite and etc.
\end{itemize}

\section*{OTHER EXPERIENCES}
\noindent
\textbf{Library Management System}  \hfill IISER Kolkata, India\\
\textit{Software Project} \hfill 2024
\begin{itemize}
    \item Developed an interactive library management system in C.
    \item Implemented book tracking and user management functionality.
\end{itemize}

\noindent
\textbf{Yuva Sangam Program (Phase 4)} \hfill Hariyan, India\\
\textit{Participant} \hfill  2024

% \section*{ENGLISH \& GRE TESTS}
%     \begin{tabularx}{1\textwidth}{
%     >{\raggedright\arraybackslash}X 
%    >{\raggedright\arraybackslash}X }
%       \textbf{IELTS/TOEFL:} \# score (overall)&   \textbf{GRE General Test:}\\ \\
%     Listening: \#XX | Reading: \#XX & Quant: \#XXX | Verbal: \#XXX\\
%     Speaking: \#XX   | Writing: \#XX& Analytical writing: \#X\\
%     Test date: \#Month \#Year &Test date: \#Month \#Year
%     \end{tabularx}


\section*{REFERENCES}
\textbf{Prof. Supriyo Mitra}\\
\textit{Faculty, Department of Earth Sciences, IISER Kolkata, Kolkata, India}\\
E-mail: \href{mailto:supriyomitra@iiserkol.ac.in}{supriyomitra@iiserkol.ac.in}\\
Scholar Profiles: \href{https://www.iiserkol.ac.in/~supriyomitra/}{Personal Page} - \href{https://scholar.google.com/citations?user=_zF1Zr8AAAAJ&hl=en}{Google Scholar} \\ \\
\textbf{Dr. Martin Vall\'ee}\\
\textit{Researcher, Institut de Physique du Globe de Paris (IPGP), Paris, France}\\
E-mail: \href{mailto:vallee@ipgp.fr}{vallee@ipgp.fr}\\
Scholar Profiles: \href{https://www.ipgp.fr/~vallee/}{Personal Page} - \href{https://scholar.google.com/citations?user=d39g-UgAAAAJ&hl=fr}{Google Scholar}  \\ \\
\textbf{Dr. Arjun Datta}\\
\textit{Assistant Professor, Department of Earth and Climate Science, Indian Institute of Science Education and Research (IISER) Pune, Pune, India}\\
E-mail: \href{mailto:arjundatta@iiserpune.ac.in }{arjundatta@iiserpune.ac.in }\\
Scholar Profiles: \href{https://sites.google.com/view/arjun-datta}{Personal Page} | \href{https://scholar.google.com/citations?hl=en&user=tgjqIncAAAAJ&view_op=list_works&sortby=pubdate}{Google Scholar} \\ \\
% End document
\end{document}